% ============================================================================
% Chapter 4: Implementation and Results  (FYDP-II)
%
% Paste after Chapter 3. Compiles with pdfLaTeX; no packages beyond a standard
% report preamble are needed, and no Bangla script is used (pdfLaTeX cannot
% typeset it without extra setup).
%
% Citation keys used here. Replace the first group with the keys your .bib or
% thebibliography already uses for these FYDP-I references:
%   lettucedetect -> FYDP-I ref [1]    luna -> [2]    ragtruth -> [16]
% The second group is new; entries are in new_references.tex:
%   banglabert, electra, lora, mmbert, gemma4, llamacpp
% ============================================================================

\chapter{Implementation and Results}
\label{ch:implementation}

This chapter describes how the framework designed in Chapter~3 was built and
evaluated. Section~\ref{sec:env} covers the hardware, software, dataset and
training configuration. Section~\ref{sec:eval} defines the evaluation protocol
and the systems the framework is compared against. Section~\ref{sec:results}
reports the results and discusses what they show.

\paragraph{Change of language from FYDP-I.} The FYDP-I design planned to train
and evaluate on the English RAGTruth corpus~\cite{ragtruth}. In FYDP-II the
primary language became Bangla. No token-level hallucination detector or
hallucination dataset exists for Bangla, which is the gap recorded in row~4 of
Table~2.1, so this is where a small, local detector adds the most value. Luna
and LettuceDetect~\cite{luna,lettucedetect} are trained and reported on English
RAGTruth only, so their published scores cannot be compared directly with ours.
Evaluating the framework on RAGTruth is left as future work.

% ----------------------------------------------------------------------------
\section{Environment Setup}
\label{sec:env}

\subsection{Hardware}

All training, evaluation and latency measurement ran on one consumer desktop
computer (Table~\ref{tab:hardware}). No cloud GPU was used, although
Section~3.1.4 had planned Google Colab for fine-tuning.

\begin{table}[htbp]
\centering
\caption{Hardware used for all experiments}
\label{tab:hardware}
\begin{tabular}{ll}
\hline
\textbf{Component} & \textbf{Specification} \\
\hline
CPU & AMD Ryzen 5 7500F, 6 cores \\
Memory & 16 GB RAM \\
GPU & AMD Radeon RX 6600, 8 GB VRAM \\
Operating system & Windows 11 Pro \\
Cloud compute & None \\
\hline
\end{tabular}
\end{table}

The RX 6600 is not supported by AMD's ROCm platform on Windows, so GPU training
used Microsoft DirectML through the \texttt{torch-directml} package. DirectML
needed two adjustments. First, it cannot compute the gradient of a layer
normalisation without a bias term, which the mmBERT comparison model uses
throughout; we re-implemented these layers with elementary tensor operations,
which give the same output to within $2\times10^{-6}$. Second, a batch of 16
sequences did not fit in 8~GB of GPU memory, so BanglaBERT was trained with a
batch of 8 and two-step gradient accumulation, which keeps the effective batch
size at 16. On this GPU, training was about four times faster than on the CPU.

\subsection{Software}

Table~\ref{tab:software} lists the software stack. Two Python environments were
kept: a CPU environment used for the demo, the API and CPU latency, and a GPU
environment with the older PyTorch release that \texttt{torch-directml}
requires. The LLM judges used for comparison (Section~\ref{sec:systems}) were
served with llama.cpp~\cite{llamacpp} through its Vulkan backend on the same GPU.

\begin{table}[htbp]
\centering
\caption{Software stack}
\label{tab:software}
\begin{tabular}{lll}
\hline
\textbf{Purpose} & \textbf{Tool} & \textbf{Version} \\
\hline
Language & Python & 3.12.8 \\
Deep learning (CPU) & PyTorch & 2.14 \\
Deep learning (GPU) & PyTorch + torch-directml & 2.4.1 + 0.2.5 \\
Models and tokenizers & Hugging Face Transformers & 4.57.6 \\
LoRA fine-tuning & Hugging Face PEFT & 0.21.0 \\
Metrics & scikit-learn & 1.9.1 \\
Data handling & pandas & 3.0.6 \\
Backend API & FastAPI + Uvicorn & 0.141.1 + 0.53.0 \\
Web demo & Streamlit & 1.64.0 \\
LLM judges & llama.cpp (Vulkan) & build b11223 \\
\hline
\end{tabular}
\end{table}

The code is organised into separate modules for data loading and validation,
the model, training, evaluation, the lexical baseline, latency measurement and
the LLM judges. A single configuration file selects the backbone model, so
every model in this chapter was trained by the same script.

\subsection{Dataset}

The framework was trained and evaluated on a Bangla RAG hallucination dataset
built by the team. Each row holds a retrieved context, a question, an answer,
a label for the whole answer (faithful or hallucinated) and a label for every
answer word. Table~\ref{tab:dataset} summarises it.

\begin{table}[htbp]
\centering
\caption{Bangla RAG hallucination dataset}
\label{tab:dataset}
\begin{tabular}{ll}
\hline
\textbf{Property} & \textbf{Value} \\
\hline
Rows & 4,000 (2,000 faithful, 2,000 hallucinated) \\
Distinct contexts & 943, each used for 2, 4 or 6 rows \\
Domains & 8, with 500 rows each \\
Context length & 71 words (median), 219 (maximum) \\
Answer length & 15 words (median), 51 (maximum) \\
Labels & one label per answer word (0 = grounded, 1 = hallucinated) \\
Sources & Bangla Wikipedia (3,658 rows), government portals (342 rows) \\
\hline
\end{tabular}
\end{table}

The eight domains are Bangladesh Affairs, History, Healthcare, Agriculture,
Government Services, Education, Universities, and Science \& Technology. The
government sources are the Directorate General of Health Services (DGHS), the
national information portal, the Ministry of Education and BANBEIS. The
hallucinated answers are spread evenly over six types: fabricated information
(334 rows), entity replacement (334), omission (333), contradiction (333),
number error (333) and date error (333).

The labels were produced by the dataset generator, not by people. The team is
carrying out human inter-annotation of the dataset; all results in this chapter
use the generated labels.

\paragraph{Preprocessing.} Three steps prepared the data.
\begin{enumerate}
  \item \textbf{Validation.} Every load checks that each hallucinated answer has
  at least one flagged word, that no faithful answer has one, and that the
  flagged words match the recorded hallucinated span. All 4,000 rows passed.
  \item \textbf{Word segmentation.} The dataset splits punctuation, such as the
  Bangla full stop (danda), into separate words but keeps numbers such as
  10,000 whole. We implemented the same rule, verified that it reproduces the
  dataset's words for all 4,000 rows, and use it again at inference, so the
  deployed model sees words exactly as it did during training.
  \item \textbf{Split by context.} Each context appears 2 to 6 times with
  different questions. A random split of rows would put the same context in
  both training and test data and inflate the test score. We therefore split
  by context, 70/15/15, keeping every domain in every part: 2,824 training rows
  (659 contexts), 588 validation rows (142 contexts) and 588 test rows (142
  contexts). No context appears in more than one part, and each part is exactly
  50\% hallucinated.
\end{enumerate}

\subsection{Model and Training Configuration}

\paragraph{Input.} Each example is encoded as
\texttt{[CLS] question [SEP] context [SEP] answer [SEP]}. The maximum length of
288 subword tokens was measured from the data (it covers 99\% of examples)
rather than fixed in advance. When an input is too long, tokens are removed from
the context, never from the answer, so no answer label is lost. Only the first
subword of each answer word carries a training label; question, context and
padding tokens are ignored by the loss.

\paragraph{Model.} The encoder is BanglaBERT~\cite{banglabert}, an ELECTRA-base
discriminator~\cite{electra} pretrained on 27.5~GB of Bangla text, with 12
layers, 12 attention heads and a hidden size of 768. On top of it sits a token
classification head: dropout of 0.1 followed by one linear layer from 768 to 2
outputs (1,538 parameters). The full model has 110,028,290 parameters, well
under the 500 million set in Section~3.1.1. BanglaBERT was chosen over
multilingual encoders because its tokenizer splits Bangla into far fewer pieces:
1.14 subword tokens per answer word, against 3.51 for mmBERT. Shorter sequences
make the model faster at both training and inference.

\paragraph{Training.} Table~\ref{tab:hyper} lists the settings. Two variants
were trained: full fine-tuning, which updates all weights, and LoRA~\cite{lora},
which freezes the encoder and trains small low-rank matrices added to the
query, key and value projections of all 12 attention layers, together with the
classification head. LoRA trains 886,274 parameters, 0.8\% of the model. Each
variant was trained three times with different random seeds.

\begin{table}[htbp]
\centering
\caption{Training settings}
\label{tab:hyper}
\begin{tabular}{ll}
\hline
\textbf{Setting} & \textbf{Value} \\
\hline
Loss & Cross-entropy over answer words \\
Optimiser & AdamW, weight decay 0.01 \\
Learning rate & $3\times10^{-5}$ (full), $3\times10^{-4}$ (LoRA) \\
Schedule & 10\% linear warm-up, then linear decay \\
Gradient clipping & 1.0 \\
Effective batch size & 16 (8 $\times$ 2 accumulation steps) \\
Epochs & up to 10, early stopping after 2 epochs without gain \\
Model selection & best validation word-level F1 \\
LoRA & rank 16, $\alpha$ = 32, dropout 0.1, on query/key/value \\
Seeds & 42, 43, 44 \\
\hline
\end{tabular}
\end{table}

\paragraph{Inference and output.} The model gives each answer word a
probability of being hallucinated. Words above 0.5 are flagged, consecutive
flagged words are merged into highlighted spans, and the answer is marked as
hallucinated when any word is flagged. Its overall score is the highest word
probability. The same checkpoint serves the Streamlit demo and a FastAPI
backend (hosted with Uvicorn) whose \texttt{/detect} endpoint returns the
verdict, the spans, the probability of every word and the time taken.

% ----------------------------------------------------------------------------
\section{Evaluation}
\label{sec:eval}

\subsection{Metrics}

The FYDP-I plan used an F1-score over whole answers, with a target above 75\%
(Section~3.1.1). Two findings led us to change the headline metrics. On the
first 2,652-row version of the dataset, a string-matching rule with no model
already reached an answer-level F1 of 0.94, above that target. And an
answer-level score does not test what the system actually shows the user, which
is the exact words that are wrong. The headline metrics are therefore measured
at word and span level:

\begin{itemize}
  \item \textbf{Word-level precision, recall and F1} for the hallucinated class,
  computed over every answer word in the test set.
  \item \textbf{Span-level F1.} A span is a run of consecutive flagged words. A
  predicted span counts as correct only if its start and end both match a gold
  span. Spans predicted on faithful answers count as false positives.
  \item \textbf{Span exact match:} the share of hallucinated answers in which
  every gold span is predicted exactly. \textbf{Span partial match:} the share
  in which some predicted span overlaps a gold span.
\end{itemize}

Answer-level precision, recall, F1, accuracy and AUROC are also reported,
always next to the baseline, together with the \textbf{false alarm rate}: the
share of faithful answers wrongly flagged. All metrics are also broken down by
hallucination type.

\paragraph{Latency.} Latency is the time to check one example at batch size 1,
measured separately on the CPU and the GPU of the machine in
Table~\ref{tab:hardware}, with nothing else running. Each measurement uses 20
warm-up calls, then 100 timed calls cycling through the same 30 test examples
(145 subword tokens on average for BanglaBERT), and reports the median and the
95th percentile.

\subsection{Systems Compared}
\label{sec:systems}

\begin{enumerate}
  \item \textbf{Lexical baseline.} Flags every answer word that does not occur
  in the context, allowing for Bangla case suffixes by prefix matching. It needs
  no training and shows how much of the task simple word overlap can solve.
  \item \textbf{BanglaBERT, full fine-tuning} (the proposed system).
  \item \textbf{BanglaBERT with LoRA} (the proposed system, trained with
  parameter-efficient fine-tuning).
  \item \textbf{mmBERT with LoRA.} mmBERT~\cite{mmbert} is the multilingual
  version of ModernBERT, the architecture FYDP-I planned to use. The original
  English ModernBERT cannot handle Bangla text well, so the multilingual model
  was used instead. It has 307 million parameters and needed a maximum length
  of 832 tokens for the same data.
  \item \textbf{LLM-as-a-judge.} Two instruction-tuned Gemma~4
  models~\cite{gemma4}, E4B (4-bit, 4.3~GB) and 12B (4-bit
  quantisation-aware, 6.4~GB), run locally with llama.cpp on the same GPU. Each
  judge received the context, question and answer, English instructions and
  two worked examples taken from the training split. It returned the
  unsupported phrases as JSON, which were mapped back onto answer words.
  Decoding was deterministic (temperature 0) with the model's reasoning mode
  switched off.
\end{enumerate}

All systems were scored on the same 588 test answers by the same evaluation
code. BanglaBERT results are the mean $\pm$ standard deviation over three
seeds; the other systems are single runs.

% ----------------------------------------------------------------------------
\section{Results and Discussion}
\label{sec:results}

\subsection{Main Results}

Table~\ref{tab:main} gives the results on the test split.

\begin{table}[htbp]
\centering
\small
\caption{Results on the test split (588 answers, 50\% hallucinated)}
\label{tab:main}
\begin{tabular}{lcccccc}
\hline
\textbf{System} & \textbf{Word P} & \textbf{Word R} & \textbf{Word F1} & \textbf{Span F1} & \textbf{Answer F1} & \textbf{False alarms} \\
\hline
Lexical baseline & 0.295 & 0.497 & 0.370 & 0.089 & 0.686 & -- \\
Gemma 4 E4B judge & 0.581 & 0.801 & 0.674 & 0.351 & 0.838 & 38\% \\
Gemma 4 12B judge & 0.753 & 0.757 & 0.755 & 0.479 & 0.849 & 26\% \\
mmBERT + LoRA & 0.817 & 0.800 & 0.809 & 0.527 & 0.894 & 10\% \\
BanglaBERT + LoRA & 0.862 & 0.786 & 0.822 & 0.513 & 0.925 & 6\% \\
\textbf{BanglaBERT (full)} & \textbf{0.883} & 0.793 & \textbf{0.835} & \textbf{0.583} & 0.924 & \textbf{4\%} \\
\hline
\end{tabular}

\smallskip
\footnotesize BanglaBERT rows are means over three seeds; standard deviations
for full fine-tuning are 0.007 (word F1) and 0.016 (span F1), and for LoRA
0.016 and 0.017. False alarms = faithful answers wrongly flagged.
\end{table}

Full fine-tuning of BanglaBERT gives the best result on every headline metric:
a word-level F1 of $0.835 \pm 0.007$ and a span-level F1 of $0.583 \pm 0.016$.
That is 0.465 above the lexical baseline in word F1, and more than six times
its span F1. The baseline's weakness is itself informative: in about a quarter
of the hallucinated answers, every wrong word also appears somewhere in the
context, so word overlap cannot detect them and the model has to use meaning.
At answer level the model reaches an F1 of 0.924 (precision 0.957, recall
0.893) and an AUROC of 0.977.

\subsection{Comparison with LLM-as-a-Judge}

Both LLM judges are less accurate than BanglaBERT. The larger judge, Gemma~4
12B, reaches a word F1 of 0.755 and a span F1 of 0.479, against 0.835 and 0.583.
The judges also raise far more false alarms: they flag 26\% (12B) and 38\% (E4B)
of faithful answers, against 4\% for BanglaBERT. In a deployed system a warning
that is wrong one time in four would soon be ignored.

The cost difference is larger still (Table~\ref{tab:efficiency}). On the same
GPU, BanglaBERT checks an answer in 11.5~ms and the 12B judge in 2,560~ms, about
220 times longer. Checking the whole test set took the 12B judge 25.7 minutes
and BanglaBERT about 7 seconds. BanglaBERT's saved model is 0.42~GB against
6.4~GB for the 12B judge. This confirms, with measurements on the same data and
hardware, the cost argument that motivated the project (Section~1.2).

The judges were stronger than BanglaBERT on one type: entity replacement
(word F1 0.876 for the 12B judge against 0.643), which suggests they use world
and linguistic knowledge that a 110-million-parameter encoder lacks. We report
this as a limitation of the proposed model and a direction for future work.

\begin{table}[htbp]
\centering
\small
\caption{Model size and speed}
\label{tab:efficiency}
\begin{tabular}{lrrrrr}
\hline
\textbf{System} & \textbf{Parameters} & \textbf{Trained} & \textbf{Saved size} & \textbf{CPU latency} & \textbf{GPU latency} \\
\hline
BanglaBERT (full) & 110.0M & 110.0M & 420 MB & 57.1 ms & 11.5 ms \\
BanglaBERT + LoRA & 110.9M & 0.89M & 3.4 MB$^{*}$ & 66.3 ms & 12.4 ms \\
mmBERT + LoRA & 308.6M & 1.08M & 4.1 MB$^{*}$ & 210.2 ms & 52.7 ms \\
Gemma 4 E4B judge & -- & -- & 4.3 GB & -- & 950 ms \\
Gemma 4 12B judge & $\sim$12B & -- & 6.4 GB & -- & 2,560 ms \\
\hline
\end{tabular}

\smallskip
\footnotesize Latency is the median time for one answer at batch size 1.
$^{*}$LoRA adapter only; it is loaded on top of the shared base model.
\end{table}

\subsection{BanglaBERT versus mmBERT}

mmBERT has 2.8 times as many parameters as BanglaBERT but was less accurate
(word F1 0.809 against 0.835) and much slower. Its multilingual tokenizer splits
Bangla into about three times as many pieces (3.51 against 1.14 subword tokens
per answer word), so its inputs are longer: a median of 383 tokens against 135.
It takes 210~ms per answer on the CPU, which breaks the 200~ms requirement,
while BanglaBERT takes 57~ms. For Bangla, a model pretrained on Bangla text is
the better choice than a larger multilingual one.

\subsection{Full Fine-Tuning versus LoRA}

LoRA trains 124 times fewer parameters than full fine-tuning (886,274 against
110 million). Its word F1 (0.822) is within the seed-to-seed variation of full
fine-tuning, so the two cannot be told apart on that metric. Its span F1 is
clearly lower (0.513 against 0.583): LoRA finds hallucinated words about as
well, but marks their boundaries less accurately. In return, LoRA trained
1.8 times faster per epoch on the GPU (168 against 304 seconds), and its saved
adapter is 3.4~MB instead of 420~MB, so one base model can serve many
domain-specific adapters. On the CPU the time saving was only 9\%, because the
backward pass still runs through the frozen encoder. LoRA is therefore a
trade-off rather than a free gain at this model size: it suits deployments
that need many small, cheaply retrained variants, while full fine-tuning gives
the most accurate spans.

\subsection{Results by Hallucination Type}

Table~\ref{tab:types} breaks word F1 down by type.

\begin{table}[htbp]
\centering
\small
\caption{Word-level F1 by hallucination type (test split)}
\label{tab:types}
\begin{tabular}{lrcccc}
\hline
\textbf{Type} & \textbf{Answers} & \textbf{Baseline} & \textbf{Gemma 12B} & \textbf{BanglaBERT} & \textbf{+ LoRA} \\
\hline
Fabricated information & 44 & 0.756 & 0.962 & \textbf{0.976} & 0.980 \\
Number error & 59 & 0.447 & 0.842 & \textbf{0.882} & 0.837 \\
Contradiction & 58 & 0.392 & 0.786 & \textbf{0.844} & 0.829 \\
Date error & 41 & 0.528 & \textbf{0.839} & 0.827 & 0.793 \\
Entity replacement & 46 & 0.198 & \textbf{0.876} & 0.643 & 0.666 \\
Omission & 46 & 0.113 & 0.366 & \textbf{0.598} & 0.598 \\
\hline
\end{tabular}
\end{table}

Fabricated information is almost always caught, because added facts have no
support anywhere in the context. Numbers, contradictions and dates are detected
well (word F1 0.83 to 0.88). The two weakest types are omission and entity
replacement. An omission leaves out part of the correct information, so the
words that remain are all supported by the context, and the model has to
notice that something is missing. An entity replacement swaps in a name that
often appears elsewhere in the same context. These are also the types the
lexical baseline almost never detects (word F1 0.11 and 0.20), so they are where
the model adds most over word overlap and where there is most room to improve.

\subsection{Training Observations}

Validation word F1 levelled off after about three epochs; the best epochs for
the three full fine-tuning runs were 3, 9 and 4, and training for up to 10
epochs gave no gain over the earlier 4-epoch run. Tuning the decision threshold
on the validation split did not help either: the chosen thresholds varied from
0.30 to 0.65 between seeds, and test word F1 was the same as with the default
of 0.5, which the deployed system therefore keeps. The standard deviation
across seeds was small (0.007 in word F1), so differences between systems
smaller than about 0.02 should not be read as real.

\subsection{Requirements Check}

Table~\ref{tab:requirements} compares the results with the requirements set in
Section~3.1.1.

\begin{table}[htbp]
\centering
\small
\caption{Requirements from FYDP-I and their outcome}
\label{tab:requirements}
\begin{tabular}{p{4.2cm}p{4.2cm}p{4.6cm}}
\hline
\textbf{Requirement} & \textbf{Target} & \textbf{Outcome} \\
\hline
Input ingestion & Context, question, answer & Met (demo and API) \\
Token-level classification & Probability per answer token & Met \\
Span aggregation & Consecutive tokens merged into spans & Met \\
Low latency & 100--200 ms per inference & Met: 57.1 ms CPU, 11.5 ms GPU \\
Lightweight deployment & At most 500M parameters, local & Met: 110M parameters, runs offline \\
Parameter-efficient training & LoRA used and measured & Met: 886K trainable parameters \\
High accuracy & F1 above 75\% on RAGTruth & Revised: word F1 0.835 on the Bangla dataset; RAGTruth not evaluated \\
User interface & Streamlit dashboard, FastAPI backend & Met \\
\hline
\end{tabular}
\end{table}

\subsection{Limitations}

\begin{itemize}
  \item \textbf{Labels not yet verified by people.} The dataset's labels were
  generated automatically. Human inter-annotation is in progress, and the
  results should be re-checked against the verified labels.
  \item \textbf{Balanced test set.} Exactly half of the test answers are
  hallucinated. Real systems hallucinate less often, which would change
  precision and the false alarm rate.
  \item \textbf{Short answers.} 96\% of answers are a single sentence. Longer
  answers with one wrong clause are harder and are not yet tested.
  \item \textbf{Context length.} BanglaBERT reads at most 512 tokens. The
  current contexts fit, but longer retrieved passages would be cut.
  \item \textbf{Scope of checking.} The system checks an answer only against
  the context it is given. It cannot detect errors when the context itself is
  wrong, and it does not fact-check text that has no source.
  \item \textbf{Comparison scope.} The LLM judges were local, 4-bit models run
  once each without reasoning. Larger cloud models might score higher, at a
  much greater cost. No comparison on English RAGTruth was made.
\end{itemize}

\subsection{Summary}

A 110-million-parameter BanglaBERT token classifier detects hallucinated words
in Bangla RAG answers with a word-level F1 of 0.835 and a span-level F1 of
0.583. It is more accurate than a lexical baseline, a larger multilingual
encoder and two local LLM judges, raises the fewest false alarms, and checks an
answer in 57~ms on a CPU or 11.5~ms on a consumer GPU. LoRA reduces the trained
parameters 124-fold and the saved model to 3.4~MB at a cost in span accuracy.
The weakest cases are omissions and entity replacements.
